\documentclass[12pt]{article}

\usepackage{setspace}
\onehalfspace

%\usepackage[utf8]{inputenc}
%\usepackage{t1enc}
%\usepackage[T1]{fontenc}

\usepackage{amsmath}

\usepackage{moodle}

\DeclareMathOperator{\Exp}{Exp}
\DeclareMathOperator{\Norm}{{\cal N}}
\DeclareMathOperator{\Binom}{Binom}
\DeclareMathOperator{\E}{E}
\DeclareMathOperator{\D}{D}
\renewcommand{\P}{\operatorname{P}}
%\DeclareMathOperator{\cov}{cov}
\DeclareMathOperator{\corr}{corr}

\DeclareMathOperator{\F}{F} 
\renewcommand{\d}{\operatorname{d\! }} %integral \! a nyero


\DeclareMathOperator{\R}{\mathbb{R}}
%\DeclareMathOperator{\Norm}{{\cal N}}
\DeclareMathOperator{\Unif}{{\cal U}}
\DeclareMathOperator{\Poi}{{Poisson}}


\DeclareMathOperator{\cov}{cov}





\begin{document}

\begin{quiz}{gyakorló}

\begin{multi}{dvv14}
$\xi$ és $\eta$ független, véges értékkészletű valségi változókra:
$$
\D^2(\xi+\eta)=\D^2(\xi)+\D^2(\eta)
$$
\item* igaz
\item hamis
\end{multi}
\begin{multi}{flenvv2}
Azt mondjuk, hogy $\xi$ és $\eta$ valségi változók függetlenek, ha
$$
\E(\xi \eta)=\E(\xi)\E(\eta)
$$
\item igaz
\item* hamis
\end{multi}
\begin{multi}{flenvv5}
Ha $\xi$ és $\eta$ diszkrét valségi változóknak létezik a szórása és nemnulla, akkor a 
korrelációjuk
$$
\corr(\xi,\eta)=\E(\xi\eta)-\E(\xi)\E(\eta)
$$
\item igaz
\item* hamis
\end{multi}
\begin{multi}{impfvv6}
Ha $\xi\sim \Unif(0,1)$ akkor  minden $a,b$-re $b\xi+a\sim \Unif(a,b)$.
\item igaz
\item* hamis
\end{multi}
\begin{multi}{felt1}
$A$-nak a $B$-re vonatkozó feltételes valsége $\P(A|B)=\frac{\P(A)}{\P(B)}$.
\item igaz
\item* hamis
\end{multi}
\begin{multi}{felt4}
$A$-nak a $B$-re vonatkozó feltételes valségét $\P(A|B)=\frac{\P(A\cdot B)}{\P(B)}$ módon értelmezzük, feltéve hogy $\P(B)>0$.
\item* igaz
\item hamis
\end{multi}
\begin{multi}{eofv10}
Bármely $\xi$-re és $a$-ra $\P(\xi= a)=0$.
\item igaz
\item* hamis
\end{multi}
\begin{multi}{impfvv9}
Ha $\xi\sim \Exp(\lambda)$ akkor a sűrűségfüggvénye 
$$
f_{\xi}(x)=
\begin{cases}
\lambda e^{-\lambda x}& \ \ x>0\\
0 & \text{máskor} 
\end{cases}
$$
\item* igaz
\item hamis
\end{multi}
\begin{multi}{sfv2}
Egy $f:\R\to [0,1]$ pontosan akkor sűrűségfüggvény, ha monoton nemcsökkenő, 
balról folytonos, $\lim_{x\to+\infty}f(x)=1$ és $\lim_{x\to-\infty}f(x)=0$
\item igaz
\item* hamis
\end{multi}
\begin{multi}{dvv16}
$\xi$ és $\eta$ véges értékkészletű valségi változókra:
$$
\D^2(\xi+\eta)=\D^2(\xi)+\D^2(\eta)
$$
\item igaz
\item* hamis
\end{multi}
\begin{multi}{dvv10}
Egy $\xi$ valségi változó az $x_1,\ldots x_n$ értékeket veheti fel. 
Ekkor a szórásnégyzete:
$$
\D^2(\xi)=\E(\xi^2)-\E(\xi)^2
$$
\item* igaz
\item hamis
\end{multi}
\begin{multi}{felt6}
$A_1,\ldots,A_n$ teljes eseményrendszer, ha $\sum A_k=\Omega$ és a tagok páronként függetlenek.
\item igaz
\item* hamis
\end{multi}
\begin{multi}{impfvv1}
Egy $\xi\sim \Unif(a,b)$ várható értéke: $\frac{a+b}{2}$
\item* igaz
\item hamis
\end{multi}
\begin{multi}{eofv7}
Bármely $\xi$-re és $a<b$-re $\P(a\le \xi <b)=F_{\xi}(b)-F_{\xi}(a)$.
\item* igaz
\item hamis
\end{multi}
\begin{multi}{dvv6}
Egy $\xi$ valségi változó az $x_1,\ldots x_n$ értékeket veheti fel. Ekkor a szórása:
$$
\D(\xi)=\sum_{k=1}^n \P(\xi=x_k)(x_k-\E(\xi))^2
$$
\item igaz
\item* hamis
\end{multi}
\begin{multi}{sfv3}
Azt mondjuk, hogy egy $f:\R\to \R$ sűrűségfüggvény 
$$
\int_{-\infty}^{+\infty}f(x)\d x\ge 0
$$
\item igaz
\item* hamis
\end{multi}
\begin{multi}{impdvv7}
Azt mondjuk, hogy $\xi\sim \Poi(\lambda)$, ha 
$$
\P(\xi=k)=e^{-\lambda}\frac{\lambda^k}{k!}\ \ \ (k=0,1,2\ldots )
$$ 
valamilyen $\lambda>0$ esetén.
\item* igaz
\item hamis
\end{multi}
\begin{multi}{esem8}
Az $A$ és $B$ események pont akkor függetlenek, ha $\P(A+B)=\P(A)+\P(B)$.
\item igaz
\item* hamis
\end{multi}
\begin{multi}{eofv4}
Minden valségi változónak van eloszlásfüggvénye.
\item* igaz
\item hamis
\end{multi}
\begin{multi}{esem9}
Ha $\P(A)=0$ akkor $A$ bekövetkezhet.
\item* igaz
\item hamis
\end{multi}
\begin{multi}{val3}
A valószínűség monoton, azaz $\P(A)<\P(B)$, ha $A\subseteq B$.
\item igaz
\item* hamis
\end{multi}
\begin{multi}{impdvv2}
Egy $\xi\sim \Binom(n,p)$ várható értéke: $n+p$
\item igaz
\item* hamis
\end{multi}
\begin{multi}{sfv9}
Ha $\xi$-nek $f$ a sűrűségfüggvénye és $F$ az eloszlásfüggvénye, akkor bármely $b$-re:
$$
F(b)=\int_{\infty}^{b}f(x)\d x
$$
\item* igaz
\item hamis
\end{multi}
\begin{multi}{esem5}
Ha $\P(A)=1$ akkor $A$ bekövetkezhet.
\item* igaz
\item hamis
\end{multi}
\begin{multi}{val4}
A valószínűség monoton, azaz $\P(A)\le \P(B)$, ha $A\subseteq B$.
\item* igaz
\item hamis
\end{multi}
\end{quiz}

\end{document}

